\documentclass{article}
\usepackage[utf8]{inputenc}
\usepackage[T1]{fontenc}
\usepackage{amsmath}
\usepackage{amssymb}
\usepackage{geometry}

\geometry{a4paper, margin=1in}

\begin{document}

\section*{Énoncé}

\subsection*{Exercice 1}

On considère dans $\mathbb{Z} \times \mathbb{Z}$ l'équation (E): $16x+19y=2017$.

\begin{enumerate}
    \item Montrer que 2017 est un nombre premier.
    \item Vérifier que le couple $(18,91)$ est une solution de (E).
    \item Résoudre dans $\mathbb{Z} \times \mathbb{Z}$ l'équation (E).
    \item Soit $(x,y)$ une solution de (E). Quelles sont les valeurs possibles de $x \wedge y$?
    \item Justifier que 637 est un inverse de 19 modulo 2017.
    \item Déterminer tous les couples $(x,y)$ solutions de (E) tels que $x \wedge y=1$.
\end{enumerate}

\end{document}
